\documentclass[10pt,a4paper]{amsart}
\usepackage[margin=19mm]{geometry}
\usepackage{amsmath,amssymb,mathtools}
\usepackage[scr=rsfs,cal=euler]{mathalfa}
\usepackage{xfrac}
\usepackage[unicode,hidelinks,bookmarks=false]{hyperref}
\newcommand{\ie}{i.e.\ }
\newcommand{\cf}{cf.\ }
\newcommand{\resp}{respectively\ }
\newcommand{\Subsection}{\subsection}
\providecommand{\nobreakdash}{\nobreak\hbox{-}}
\setlength{\emergencystretch}{2em}
\multlinegap0pt
\usepackage{bm}
\usepackage{bbm}
\usepackage{dsfont}
\usepackage{xspace}
\usepackage{cancel}
\usepackage{mathtools}
\usepackage{amsthm}
\makeatletter
\@ifundefined{defn}{\newtheorem{defn}{Defn}}{}
\@ifundefined{Cor}{\newtheorem{Cor}[defn]{Corollary}}{}
\@ifundefined{Prop}{\newtheorem{Prop}[defn]{Proposition}}{}
\makeatother
\newcommand{\Z}{\ensuremath{\mathbb{Z}}}
\title[On torsion in (bi)linearized Legendrian contact homology in ]{On torsion in (bi)linearized Legendrian contact homology in dimension 3}
\author{Fr\'ed\'eric Bourgeois, Salammbo Connolly}
\date{Source: arXiv:2510.22360v1 (CC BY 4.0)}
\begin{document}
\maketitle

\noindent\emph{Source and licence.} On torsion in (bi)linearized Legendrian contact homology in dimension 3. Version arXiv:2510.22360v1 (CC BY 4.0), distributed under the Creative Commons Attribution 4.0 International licence, as recorded in the arXiv licence metadata for this exact version and shown on the abstract page https://arxiv.org/abs/2510.22360v1. Full rights notice: licenses/arxiv-2510.22360-CC-BY-4.0.txt.\par\smallskip

\noindent\textit{Editorial scope.} Counted unit: the Cor whose source label is \texttt{cor:torsion\_iso} in the article (source0.tex lines 302--303, source label \texttt{cor:torsion\_iso}), delivered together with the article's own complete original proof of it (source0.tex lines 305--320, one \texttt{proof} environment of 16 lines). No mathematics is written, repaired or re-derived: statement and proof are the article's own text line for line. Required context retained uncounted: eq:dualityseq-bilin (lines 258--262); prop:tau (lines 266--272). Every definition, display and label of those lines is retained unchanged. Omitted from this package and not redistributed: 1--257, 263--265, 273--301, 304--304, 321--1208. Nothing inside a retained excerpt is omitted or altered: every retained body is delivered exactly as the article's lines are written, so no omission receipt of any kind accompanies this package. Citations: the retained excerpts carry no citation command at all, so no bibliography entry of the article is transcribed and no reference is redistributed. Reference adaptation: none was needed and none was made; every reference inside the delivered unit resolves inside it, so no unresolved reference or citation is produced. Printed slips kept literal: no printed word, symbol, hypothesis or number of the counted statement or of its proof was altered. Layout and class: the article's own class is \texttt{amsart} with packages including graphics, epsfig, psfrag, fontenc, inputenc, amsmath, amsfonts, latexsym, amssymb, color, amsthm, pinlabel, textcomp, subcaption; the wrapper is the lane's pinned \texttt{amsart} editorial wrapper at 10pt on A4, which restates the theorem-like environments the retained text needs and adds the article's own macro definitions that the retained text needs (\texttt{\textbackslash Z}), with extra wrapper packages (bm, bbm, dsfont, xspace, cancel, mathtools). The delivered numbering is the unit's own. Assets: no retained span contains a figure environment, a graphics-inclusion command, a PDF-page inclusion, a table or a TikZ picture; no image file is copied into this package, the package declares no asset dependency and no image file is redistributed with it. Licence: version arXiv:2510.22360v1 of this article is distributed under the Creative Commons Attribution 4.0 International licence, as recorded in the arXiv licence metadata for this exact version and shown on the abstract page https://arxiv.org/abs/2510.22360v1. A full rights notice is delivered at licenses/arxiv-2510.22360-CC-BY-4.0.txt. Characters: every retained excerpt is pure ASCII, so the delivered engine, XeLaTeX with its default Latin Modern text font, reports no missing character.
\begin{equation} \label{eq:dualityseq-bilin}
\ldots \to LCH^{n-k-1}_{\varepsilon_2, \varepsilon_1}(\Lambda) \to LCH^{\varepsilon_1,\varepsilon_2}_k(\Lambda) 
\stackrel{\tau_k}{\longrightarrow} H_k(\Lambda) \stackrel{\sigma_{n-k}}{\longrightarrow} LCH^{n-k}_{\varepsilon_2, \varepsilon_1}
(\Lambda) \to \ldots
\end{equation}

\begin{Prop}  \label{prop:tau}
Let $\Lambda$ be a Legendrian knot. If $\varepsilon$ is an augmentation of $\Lambda$, then
$\tau_0: LCH_0^\varepsilon(\Lambda) \to H_0(\Lambda)$ vanishes and $\tau_1: LCH_1^\varepsilon(\Lambda) \to H_1(\Lambda)$
is surjective. If $\varepsilon_1$ and $\varepsilon_2$ are augmentations of $\Lambda$ that are not dga-homotopic, then the map
$\tau_1: LCH_1^{\varepsilon_1, \varepsilon_2}(\Lambda) \to H_1(\Lambda)$ vanishes and the map 
$\tau_0: LCH_0^{\varepsilon_1, \varepsilon_2}(\Lambda) \to H_0(\Lambda)$ does not vanish.
\end{Prop}

\begin{Cor}\label{cor:torsion_iso}Let $\Lambda$ be a Legendrian knot, and let $\varepsilon_1,\varepsilon_2$ be two $\Z$-valued augmentations of its DGA. We denote the torsion part of $LCH^{\varepsilon_1,\varepsilon_2}_k(\Lambda)$ as $T_k(\Lambda,\varepsilon_1,\varepsilon_2)$. Then if either $\varepsilon_1\sim\varepsilon_2$, or $\tau_0$ is surjective, then $\forall k\in\Z$, $T_k(\Lambda, \varepsilon_1,\varepsilon_2)\cong T_{-k-1}(\Lambda, \varepsilon_2,\varepsilon_1)$. 
\end{Cor}

\begin{proof} The fact that for any $k\neq-1,0,1$, $LCH^{\varepsilon_1,\varepsilon_2}_k(\Lambda)\cong LCH^{-k}_{\varepsilon_2,\varepsilon_1}(\Lambda)$ is a direct consequence of the long exact sequence \eqref{eq:dualityseq-bilin}, and the fact that $H_k(\Lambda)=0$ for $k\neq0,1$. Since for any $k$, by the Universal Coefficient Theorem, the torsion part of $LCH^{k}_{\varepsilon_1,\varepsilon_2}(\Lambda)$ is isomorphic to the torsion part of $LCH_{k-1}^{\varepsilon_1,\varepsilon_2}(\Lambda)$, this then gives us $T_k(\Lambda, \varepsilon_1,\varepsilon_2)\cong T_{-k-1}(\Lambda, \varepsilon_2,\varepsilon_1)$ for any $k\neq-1,0$ (the case where $k=1$ is paired with the $k=-2$ case).

Suppose now that $\Lambda$ is connected, and that $\varepsilon_1$ and $\varepsilon_2$ are dga-homotopic. We may then write the homology as $LCH_k^{\varepsilon_1}$, as it is isomorphic to the linearized homology. We want to show that $T_0(\Lambda,\varepsilon_1,\varepsilon_2)\cong T_{-1}(\Lambda,\varepsilon_1,\varepsilon_2)$ (in this case the orders of the augmentations do not matter). The long exact sequence reads
 \begin{align*}
\lefteqn{\ldots \rightarrow LCH_{1}^{\varepsilon_1}(\Lambda) \overset{\tau_1}\rightarrow H_1(\Lambda) = \Z \overset{\psi_0}\rightarrow LCH^{0}_{\varepsilon_1}(\Lambda)} \\
& \hspace{3cm} \overset{\phi_0}\rightarrow LCH_{0}^{\varepsilon_1}(\Lambda) \newline \overset{\tau_0}\rightarrow H_{0}(\Lambda) = \Z\rightarrow \ldots
\end{align*}
Proposition \ref{prop:tau} tells us that $\tau_0$ vanishes, whereas $\tau_1$ is surjective. Since $\tau_0=0$, $\phi_0$ is surjective, and since $\tau_1$ is surjective, $\psi_0=0$, and so $\phi_0$ is injective. And so $LCH_{0}^{\varepsilon_1}(\Lambda)\cong LCH^{0}_{\varepsilon_1}(\Lambda)$. In particular, in this case $T_0(\Lambda, \varepsilon_1,\varepsilon_2)$ is isomorphic to the torsion part of $LCH^{0}_{\varepsilon_1}(\Lambda)$, which is isomorphic to $T_{-1}(\Lambda,\varepsilon_1,\varepsilon_2)$. 

Let us now consider the case where $\varepsilon_1$ and $\varepsilon_2$ are not dga-homotopic. Then Proposition \ref{prop:tau} tells us that $\tau_1$ vanishes, which gives us $LCH^{\varepsilon_1,\varepsilon_2}_1(\Lambda)\cong LCH^{-1}_{\varepsilon_2,\varepsilon_1}(\Lambda)$. Furthermore, the long exact sequence reads 
 \begin{align*}
\lefteqn{\ldots \overset{\tau_1}\rightarrow H_1(\Lambda) = \Z \overset{\psi_0}\rightarrow LCH^{0}_{\varepsilon_2,\varepsilon_1}(\Lambda)  \overset{\phi_0}\rightarrow LCH_{0}^{\varepsilon_1,\varepsilon_2}(\Lambda)} \\
& \hspace{3cm} \overset{\tau_0}\rightarrow H_{0}(\Lambda) = \Z\overset{\psi_{-1}}\rightarrow LCH^{1}_{\varepsilon_2,\varepsilon_1}(\Lambda)  \overset{\phi_{-1}}\rightarrow LCH_{-1}^{\varepsilon_1,\varepsilon_2}(\Lambda)\rightarrow0.
\end{align*}
Then if $\tau_0$ is surjective, then $\psi_{-1}$ must vanish, and so $\psi_{-1}$ must be injective. It is also clearly surjective due to $H_{-1}(\Lambda)$ being trivial, and so we have an isomorphism. $T_{-1}(\Lambda,\varepsilon_1,\varepsilon_2)$ must then be isomorphic to the torsion part of $LCH^{1}_{\varepsilon_2,\varepsilon_1}(\Lambda) $, which is isomorphic to $T_{0}(\Lambda,\varepsilon_2,\varepsilon_1)$.
\end{proof}

\end{document}
